\section{Caos}

Engenharia do caos é experimento, não sabotagem. O roteiro:

\begin{enumerate}
\item definir o estado estável;
\item formular uma hipótese;
\item injetar uma falha real;
\item tentar refutar a hipótese.
\end{enumerate}

Comece pequeno, limitando o \termo{raio de explosão}{\emph{Blast radius}:
o tanto do sistema e dos usuários que um experimento pode afetar.}, e
tenha uma condição de parada.

O caos põe à prova os mecanismos de resiliência: \emph{timeout}, \emph{retry}
com \emph{backoff}, \emph{circuit breaker} (\citetitle{artigo:circuitbreaker}),
\emph{bulkhead} e \emph{fallback}. Para injetar falhas, a aula apresentou
\citetitle{doc:chaos-monkey-spring} e \citetitle{doc:toxiproxy}; o
\citetitle{doc:netflix-chaosmonkey} é a origem da ideia, e
\citetitle{doc:resilience4j} cobre os padrões de resiliência em Java.
